\documentclass[12pt,a4paper]{article}
\usepackage[UTF8]{ctex}
\usepackage{graphicx}
\usepackage{geometry}
\usepackage{ulem}
\usepackage{enumitem}
\usepackage{titlesec}
\usepackage{multicol}
\usepackage[export]{adjustbox}
\usepackage{amsmath}
\usepackage[utf8]{inputenc}
\usepackage{CJKutf8}
\usepackage{fancyhdr}
\usepackage{lastpage}
\usepackage{ulem}
\usepackage{tabularx}
\usepackage{float}
\usepackage{siunitx}
\usepackage{multirow} 
\usepackage{booktabs}
\usepackage{pgfplots}
\pgfplotsset{compat=1.18}
\usepackage{tikz}

% 页眉页脚设置
\pagestyle{fancy}
\fancyhf{}
\fancyhead[C]{\small
    \makebox[4.5em][c]{实验名称：} \uline{\makebox[17em][c]{喇叭天线的辐射特性测量}} 
    \makebox[2.5em][c]{姓名：} \uline{\makebox[4em][c]{张赫}} 
    \makebox[2.5em][c]{学号：} \uline{\makebox[5em][c]{3240101459}}
}
\fancyfoot[C]{\small 第 \thepage\ 页，共 \pageref{LastPage} 页}

\renewcommand{\headrulewidth}{0pt}
\renewcommand{\footrulewidth}{0pt}
\setlength{\headheight}{50pt}

\fancypagestyle{plain}{
    \fancyhf{}
    \renewcommand{\headrulewidth}{0pt}
    \fancyfoot[C]{第 \thepage\ 页，共 \pageref{LastPage} 页}
}

\geometry{left=2.54cm,right=1.5cm,top=2.54cm,bottom=2.54cm}
\renewcommand{\normalsize}{\fontsize{12pt}{16pt}\selectfont}
\titleformat{\section}{\fontsize{16pt}{22pt}\selectfont\bfseries}{\thesection}{1em}{}
\titleformat{\subsection}{\fontsize{15pt}{20pt}\selectfont\bfseries}{\thesubsection}{1em}{}

\begin{document}
\thispagestyle{plain}

\begin{multicols}{2}
    \noindent
    ~\\
    ~\\
    ~\\
    \hphantom{占位}
    \raisebox{-0.5\baselineskip}
    {
\includegraphics[width=1.85in,height=0.58in]{ZJU.jpg}}
    \hspace{.5em}
    \raisebox{0pt}[\height][0pt]
    {\textbf{\fontsize{26pt}{30pt}\selectfont 实验报告}}
    
    \columnbreak
    
    \begin{flushright}
        \small
        \makebox[3em][l]{专业：} \uline{\makebox[9em][c]{法语-电子科学与技术}} \\
        \makebox[3em][l]{姓名：} \uline{\makebox[9em][c]{张赫}} \\
        \makebox[3em][l]{学号：} \uline{\makebox[9em][c]{3240101459}} \\
        \makebox[3em][l]{日期：} \uline{\makebox[9em][c]{2026年6月17日}} \\
        \makebox[3em][l]{地点：} \uline{\makebox[9em][c]{紫金港东四221}} \\
    \end{flushright}   
\end{multicols}

{\small
\noindent
课程名称： \uline{电磁场与电磁波} \hfill 指导老师： \uline{李凯、王子立} \hfill 成绩： \uline{\hspace{3cm}} \\
实验名称：\uline{喇叭天线的辐射特性测量}\hfill 实验类型：\uline{验证实验}\hfill 同组学生姓名：\uline{李雨泽、赵奕铭}
}

\section*{一、实验目的}
\begin{enumerate}
    \item 掌握天线S参数、驻波比(VSWR)、反射系数($\Gamma$)、回波损耗(RL)等电路特性参数的测量方法及其相互转换关系。
    \item 理解天线辐射方向图的概念，掌握E面和H面方向图的测量方法。
    \item 认识天线的极化特性，理解极化匹配与交叉极化隔离度的概念。
    \item 验证电磁波在空间传播中功率与距离的关系（弗里斯传输公式）。
\end{enumerate}

\section*{二、实验原理}
天线是无线通信系统的核心部件，负责导行波与空间电磁波的双向转换。描述天线的参量主要分为电路特性参量（输入阻抗、驻波比、反射系数、回波损耗等）和辐射特性参量（方向图、增益、极化等）。

根据观察点与天线的距离$r$不同，场区可分为感应近场区（$r < \lambda/2\pi$）、辐射近场区和辐射远场区（$r \gg 2D_HD_E/\lambda$）。远场区场为平面波，方向图稳定，是工程测量常用区域。

反射系数$\Gamma$、电压驻波比VSWR和回波损耗RL的关系为：
\[
\Gamma = \frac{\text{VSWR} - 1}{\text{VSWR} + 1}, \quad RL = -20\lg|\Gamma|
\]

弗里斯传输公式描述了电磁波在空间传播中功率与距离的关系：
\[
\frac{P_r}{P_t} = \left(\frac{\lambda}{4\pi R}\right)^2 G_t G_r
\]

角锥喇叭天线增益估算：
\[
G = 0.51\frac{4\pi A_p}{\lambda^2} \quad (A_p = D_H \times D_E)
\]

半功率波束宽度估算：
\[
\text{H面：}2\theta_{0.5H} \approx 1.18\frac{\lambda}{D_H}, \quad
\text{E面：}2\theta_{0.5E} \approx 0.89\frac{\lambda}{D_E}
\]

\section*{三、实验装置与仪器}
本实验使用天线综合测控系统，由硬件分系统和软件分系统构成。硬件分系统包括测量暗室子系统、采样架子系统和信号链路子系统；软件分系统由测量控制、数据采集处理和结果显示输出子系统组成。

\begin{table}[H]
\centering
\caption{角锥喇叭天线尺寸}
\begin{tabular}{c|c|c}
\hline
\textbf{尺寸} & \textbf{发射天线} & \textbf{接收天线} \\
\hline
$D_E$ (cm) & 3.7 & 10.5 \\
$D_H$ (cm) & 8.2 & 14.1 \\
\hline
\end{tabular}
\end{table}

主要仪器：矢量网络分析仪、角锥喇叭天线（收发各一只）、天线采样架、微波屏蔽板、极化栅网、计算机及“天线测量”程序。

\section*{四、实验步骤与数据记录}

\subsection*{1. 电路特性参数测量}

在[调试]窗口中调节收发天线极化方向与轴向角度一致后，选择[S参数测量] → [驻波比] → [S11]，设定扫描频率范围为3-8.5 GHz，测得$S_{11}$曲线。同时选择[S参数测量] → [对数幅度] → [S21]，设定扫描频率范围为6-8.5 GHz，测得$S_{21}$曲线。

\begin{figure}[H]
\centering
\begin{tikzpicture}
\begin{axis}[
    width=0.8\textwidth,
    height=6cm,
    xlabel={频率 (GHz)},
    ylabel={$S_{11}$ (dB)},
    xmin=3, xmax=8.5,
    ymin=0, ymax=30,
    grid=both,
    grid style={line width=.1pt, draw=gray!30},
    legend style={at={(0.98,0.98)},anchor=north east},
    title={图1\quad $S_{11}$回波损耗测试曲线}
]
\addplot[blue,thick,smooth] coordinates {
(3.000,5.228)(3.028,24.224)(3.055,4.516)(3.083,23.149)(3.110,5.582)
(3.138,11.138)(3.165,5.649)(3.193,12.381)(3.220,7.811)(3.248,6.891)
(3.275,9.629)(3.303,8.804)(3.330,12.192)(3.358,5.287)(3.385,11.815)
(3.413,5.597)(3.440,9.035)(3.468,4.695)(3.495,10.893)(3.523,5.685)
(3.550,8.361)(3.578,6.542)(3.605,10.461)(3.633,7.072)(3.660,6.230)
(3.688,7.320)(3.715,5.949)(3.743,6.968)(3.770,4.973)(3.798,8.811)
(3.825,5.266)(3.853,8.261)(3.880,5.279)(3.908,8.933)(3.935,4.793)
(3.963,6.420)(3.990,4.603)(4.018,6.681)(4.045,4.746)(4.073,6.513)
(4.100,5.272)(4.128,7.088)(4.155,5.780)(4.183,6.535)(4.210,5.609)
(4.238,5.668)(4.265,5.520)(4.293,4.929)(4.320,5.444)(4.348,5.020)
(4.375,6.002)(4.403,5.117)(4.430,6.511)(4.458,5.021)(4.485,5.885)
(4.513,4.555)(4.540,5.589)(4.568,4.568)(4.595,5.419)(4.623,4.715)
(4.650,6.193)(4.678,5.444)(4.705,5.872)(4.733,5.381)(4.760,5.481)
(4.788,5.007)(4.815,4.317)(4.843,4.918)(4.870,4.281)(4.898,5.255)
(4.925,4.136)(4.953,6.692)(4.980,4.641)(5.008,6.765)(5.035,4.468)
(5.063,7.406)(5.090,4.315)(5.118,5.598)(5.145,4.107)(5.173,5.888)
(5.200,4.545)(5.228,5.235)(5.255,5.478)(5.283,5.466)(5.310,5.591)
(5.338,4.459)(5.365,6.786)(5.393,3.959)(5.420,6.164)(5.448,3.813)
(5.475,10.613)(5.503,3.828)(5.530,10.137)(5.558,4.474)(5.585,15.041)
(5.613,4.426)(5.640,9.014)(5.668,5.451)(5.695,6.289)(5.723,4.582)
(5.750,5.438)(5.778,6.854)(5.805,3.796)(5.833,5.307)(5.860,4.170)
(5.888,7.232)(5.915,3.070)(5.943,5.762)(5.970,3.796)(5.998,6.169)
(6.025,3.241)(6.053,8.259)(6.080,6.063)(6.108,6.793)(6.135,5.210)
(6.163,7.902)(6.190,6.755)(6.218,3.305)(6.245,4.035)(6.273,3.139)
(6.300,4.371)(6.328,2.638)(6.355,6.699)(6.383,3.313)(6.410,4.898)
(6.438,3.066)(6.465,6.645)(6.493,3.844)(6.520,3.480)(6.548,3.734)
(6.575,4.333)(6.603,3.328)(6.630,3.036)(6.658,4.177)(6.685,3.283)
(6.713,2.754)(6.740,2.700)(6.768,3.012)(6.795,2.093)(6.823,1.975)
(6.850,2.203)(6.878,2.194)(6.905,1.678)(6.933,1.877)(6.960,2.234)
(6.988,1.681)(7.015,1.384)(7.043,1.638)(7.070,1.832)(7.098,1.237)
(7.125,1.188)(7.153,1.325)(7.180,1.409)(7.208,1.295)(7.235,1.177)
(7.263,1.549)(7.290,1.435)(7.318,1.290)(7.345,1.216)(7.373,1.684)
(7.400,1.415)(7.428,1.206)(7.455,1.393)(7.483,1.775)(7.510,1.422)
(7.538,1.425)(7.565,1.728)(7.593,1.517)(7.620,1.163)(7.648,1.443)
(7.675,1.567)(7.703,1.152)(7.730,1.094)(7.758,1.347)(7.785,1.386)
(7.813,1.076)(7.840,1.198)(7.868,1.259)(7.895,1.273)(7.923,1.236)
(7.950,1.297)(7.978,1.124)(8.005,1.137)(8.033,1.320)(8.060,1.468)
(8.088,1.471)(8.115,1.161)(8.143,1.082)(8.170,1.232)(8.198,1.439)
(8.225,1.305)(8.253,1.151)(8.280,1.044)(8.308,1.305)(8.335,1.409)
(8.363,1.442)(8.390,1.252)(8.418,1.057)(8.445,1.338)(8.473,1.484)
(8.500,1.453)
};
\end{axis}
\end{tikzpicture}
\end{figure}

\begin{figure}[H]
\centering
\begin{tikzpicture}
\begin{axis}[
    width=0.8\textwidth,
    height=6cm,
    xlabel={频率 (MHz)},
    ylabel={$S_{21}$ (dB)},
    xmin=6000, xmax=8500,
    ymin=-75, ymax=-25,
    grid=both,
    grid style={line width=.1pt, draw=gray!30},
    legend style={at={(0.02,0.02)},anchor=south west},
    title={图2\quad $S_{21}$扫频测试曲线}
]
\addplot[red,thick,smooth] coordinates {
(6000.00,-69.988)(6012.50,-69.603)(6025.00,-68.437)(6037.50,-68.888)(6050.00,-68.564)
(6062.50,-67.142)(6075.00,-65.583)(6087.50,-65.077)(6100.00,-64.857)(6112.50,-61.936)
(6125.00,-60.782)(6137.50,-62.153)(6150.00,-62.370)(6162.50,-61.528)(6175.00,-60.796)
(6187.50,-60.846)(6200.00,-60.157)(6212.50,-58.604)(6225.00,-58.656)(6237.50,-59.361)
(6250.00,-59.115)(6262.50,-58.585)(6275.00,-58.176)(6287.50,-57.594)(6300.00,-56.287)
(6312.50,-55.040)(6325.00,-54.971)(6337.50,-55.282)(6350.00,-54.532)(6362.50,-53.230)
(6375.00,-52.839)(6387.50,-52.856)(6400.00,-51.608)(6412.50,-49.633)(6425.00,-48.931)
(6437.50,-49.731)(6450.00,-49.666)(6462.50,-47.720)(6475.00,-46.678)(6487.50,-47.612)
(6500.00,-46.739)(6512.50,-44.267)(6525.00,-43.309)(6537.50,-44.646)(6550.00,-44.714)
(6562.50,-43.239)(6575.00,-42.490)(6587.50,-42.412)(6600.00,-41.421)(6612.50,-39.967)
(6625.00,-39.432)(6637.50,-39.688)(6650.00,-39.587)(6662.50,-39.203)(6675.00,-38.222)
(6687.50,-37.401)(6700.00,-36.534)(6712.50,-35.318)(6725.00,-34.637)(6737.50,-34.582)
(6750.00,-35.143)(6762.50,-35.120)(6775.00,-34.614)(6787.50,-34.448)(6800.00,-34.395)
(6812.50,-33.673)(6825.00,-32.937)(6837.50,-33.415)(6850.00,-33.944)(6862.50,-33.267)
(6875.00,-32.094)(6887.50,-31.892)(6900.00,-31.919)(6912.50,-31.655)(6925.00,-30.801)
(6937.50,-30.454)(6950.00,-30.551)(6962.50,-30.193)(6975.00,-29.235)(6987.50,-28.590)
(7000.00,-28.509)(7012.50,-28.041)(7025.00,-27.639)(7037.50,-27.347)(7050.00,-27.534)
(7062.50,-27.597)(7075.00,-27.297)(7087.50,-27.228)(7100.00,-27.505)(7112.50,-27.476)
(7125.00,-27.193)(7137.50,-27.240)(7150.00,-27.702)(7162.50,-28.027)(7175.00,-27.864)
(7187.50,-27.673)(7200.00,-27.976)(7212.50,-28.144)(7225.00,-27.757)(7237.50,-27.700)
(7250.00,-28.194)(7262.50,-28.609)(7275.00,-28.422)(7287.50,-28.159)(7300.00,-28.295)
(7312.50,-28.439)(7325.00,-28.301)(7337.50,-28.267)(7350.00,-28.715)(7362.50,-29.190)
(7375.00,-29.289)(7387.50,-29.128)(7400.00,-29.023)(7412.50,-28.976)(7425.00,-28.926)
(7437.50,-28.787)(7450.00,-28.776)(7462.50,-28.822)(7475.00,-28.689)(7487.50,-28.357)
(7500.00,-27.673)(7512.50,-27.450)(7525.00,-27.306)(7537.50,-27.128)(7550.00,-26.938)
(7562.50,-26.839)(7575.00,-26.842)(7587.50,-26.852)(7600.00,-27.045)(7612.50,-27.148)
(7625.00,-27.180)(7637.50,-27.513)(7650.00,-27.995)(7662.50,-28.332)(7675.00,-28.614)
(7687.50,-28.811)(7700.00,-29.046)(7712.50,-29.173)(7725.00,-29.275)(7737.50,-29.122)
(7750.00,-29.287)(7762.50,-29.328)(7775.00,-29.253)(7787.50,-29.158)(7800.00,-28.891)
(7812.50,-28.568)(7825.00,-28.342)(7837.50,-28.286)(7850.00,-28.205)(7862.50,-28.250)
(7875.00,-28.306)(7887.50,-28.452)(7900.00,-28.401)(7912.50,-28.287)(7925.00,-28.303)
(7937.50,-28.552)(7950.00,-28.777)(7962.50,-28.829)(7975.00,-28.833)(7987.50,-28.683)
(8000.00,-28.589)(8012.50,-28.402)(8025.00,-28.306)(8037.50,-28.296)(8050.00,-28.277)
(8062.50,-28.350)(8075.00,-28.186)(8087.50,-27.841)(8100.00,-27.720)(8112.50,-27.688)
(8125.00,-27.639)(8137.50,-27.883)(8150.00,-28.064)(8162.50,-28.271)(8175.00,-28.465)
(8187.50,-28.825)(8200.00,-29.065)(8212.50,-29.229)(8225.00,-29.437)(8237.50,-29.598)
(8250.00,-29.635)(8262.50,-29.687)(8275.00,-29.588)(8287.50,-29.650)(8300.00,-29.704)
(8312.50,-29.717)(8325.00,-29.572)(8337.50,-29.269)(8350.00,-29.019)(8362.50,-28.643)
(8375.00,-28.302)(8387.50,-28.023)(8400.00,-27.857)(8412.50,-27.757)(8425.00,-27.677)
(8437.50,-27.555)(8450.00,-27.576)(8462.50,-27.673)(8475.00,-27.729)(8487.50,-27.801)
(8500.00,-27.834)
};
\end{axis}
\end{tikzpicture}
\end{figure}

由$S_{11}$曲线可知，在7.8-8.5 GHz频段内回波损耗较低（$S_{11} < 2$ dB），天线在该频段具有较好的阻抗匹配特性。由$S_{21}$曲线可知，在6.7-7.8 GHz范围内传输系数较高（$S_{21} > -30$ dB），在7.5 GHz附近达到最佳传输状态。

\subsection*{2. 辐射方向图测量}

设定测试频率分别为5 GHz、7 GHz和8 GHz，收发天线间距$R = 50$ cm。

\begin{figure}[H]
\centering
\begin{tikzpicture}
\begin{axis}[
    width=0.75\textwidth,
    height=6cm,
    xlabel={角度 (度)},
    ylabel={归一化幅度 (dB)},
    xmin=-135, xmax=135,
    ymin=-50, ymax=0,
    grid=both,
    grid style={line width=.1pt, draw=gray!30},
    legend style={at={(0.02,0.02)},anchor=south west},
    title={图3\quad H面辐射方向图}
]
\addplot[blue,thick] coordinates {
(-135,-3.36)(-133,-2.44)(-131,-0.72)(-129,-5.46)(-127,-4.80)(-125,-5.89)(-123,-3.27)
(-121,-4.01)(-119,-2.54)(-117,-1.50)(-115,-4.46)(-113,-4.10)(-111,-5.09)(-109,-4.83)
(-107,-2.20)(-105,-3.72)(-103,-3.14)(-101,-5.75)(-99,-2.51)(-97,-3.27)(-95,-7.00)
(-93,-5.19)(-91,-1.54)(-89,-3.25)(-87,-4.56)(-85,-2.01)(-83,-5.10)(-81,-2.07)
(-79,-6.27)(-77,-3.77)(-75,-3.92)(-73,-3.33)(-71,-5.36)(-69,-3.25)(-67,-5.66)
(-65,-4.09)(-63,-3.06)(-61,-3.21)(-59,-2.86)(-57,-3.39)(-55,-3.50)(-53,-3.03)
(-51,-4.67)(-49,-2.93)(-47,-4.75)(-45,-1.68)(-43,-3.38)(-41,-7.85)(-39,-4.29)
(-37,-1.83)(-35,-5.58)(-33,-4.08)(-31,-1.00)(-29,-3.40)(-27,-2.41)(-25,-2.37)
(-23,-6.81)(-21,-2.95)(-19,-1.81)(-17,-4.44)(-15,-7.23)(-13,-5.99)(-11,-4.89)
(-9,-1.42)(-7,-4.47)(-5,-6.15)(-3,-8.17)(-1,-4.64)(1,-4.72)(3,-4.64)(5,-3.15)
(7,-6.70)(9,-1.36)(11,-1.33)(13,-4.11)(15,-5.16)(17,-1.80)(19,-3.68)(21,-2.10)
(23,-4.63)(25,-2.76)(27,-6.22)(29,-6.27)(31,-2.35)(33,-2.78)(35,-3.09)(37,-0.65)
(39,-3.71)(41,-5.23)(43,-0.69)(45,-2.85)(47,-1.15)(49,-3.83)(51,-1.77)(53,-2.45)
(55,-4.74)(57,-3.14)(59,-1.77)(61,-4.11)(63,-6.47)(65,-6.03)(67,-4.71)(69,-4.80)
(71,-1.66)(73,-4.01)(75,-3.24)(77,-4.95)(79,-4.58)(81,-4.44)(83,-2.21)(85,-3.66)
(87,-5.50)(89,-3.17)(91,-4.85)(93,-7.49)(95,-4.31)(97,-1.76)(99,-4.27)(101,-3.47)
(103,-4.15)(105,-4.81)(107,-4.08)(109,-2.69)(111,-6.15)(113,-2.27)(115,-4.42)
(117,-6.20)(119,-3.89)(121,-0.68)(123,-2.90)(125,-1.93)(127,-6.05)(129,-4.11)
(131,-4.56)(133,-5.55)(135,0.00)
};
\addplot[red,thick] coordinates {
(-135,-34.52)(-133,-35.90)(-131,-40.15)(-129,-42.83)(-127,-38.42)(-125,-39.95)(-123,-33.35)
(-121,-30.62)(-119,-28.88)(-117,-28.41)(-115,-29.96)(-113,-29.08)(-111,-29.80)(-109,-31.32)
(-107,-33.61)(-105,-39.67)(-103,-48.30)(-101,-59.12)(-99,-34.91)(-97,-32.40)(-95,-33.40)
(-93,-28.74)(-91,-26.33)(-89,-27.24)(-87,-26.28)(-85,-26.47)(-83,-24.76)(-81,-24.35)
(-79,-26.81)(-77,-26.72)(-75,-29.66)(-73,-28.93)(-71,-28.58)(-69,-28.42)(-67,-28.25)
(-65,-27.65)(-63,-26.87)(-61,-25.98)(-59,-25.80)(-57,-25.23)(-55,-23.49)(-53,-23.12)
(-51,-22.71)(-49,-21.82)(-47,-20.98)(-45,-19.54)(-43,-18.18)(-41,-16.83)(-39,-14.74)
(-37,-14.06)(-35,-12.45)(-33,-10.70)(-31,-9.42)(-29,-8.38)(-27,-7.65)(-25,-6.10)
(-23,-5.27)(-21,-4.28)(-19,-3.65)(-17,-2.58)(-15,-2.04)(-13,-1.47)(-11,-1.14)
(-9,-0.84)(-7,-0.42)(-5,-0.20)(-3,-0.16)(-1,-0.02)(1,0.00)(3,-0.05)(5,-0.10)
(7,-0.40)(9,-0.59)(11,-0.86)(13,-1.25)(15,-1.65)(17,-2.22)(19,-2.76)(21,-3.91)
(23,-4.70)(25,-5.64)(27,-6.41)(29,-7.12)(31,-8.18)(33,-9.21)(35,-10.48)(37,-11.45)
(39,-12.77)(41,-14.34)(43,-14.88)(45,-16.05)(47,-17.79)(49,-18.51)(51,-19.00)
(53,-19.68)(55,-20.51)(57,-21.30)(59,-21.73)(61,-22.06)(63,-22.76)(65,-23.60)
(67,-24.02)(69,-24.86)(71,-25.61)(73,-27.36)(75,-28.32)(77,-28.79)(79,-27.99)
(81,-27.13)(83,-26.11)(85,-26.80)(87,-27.13)(89,-27.51)(91,-28.43)(93,-29.56)
(95,-30.63)(97,-31.10)(99,-32.32)(101,-33.98)(103,-34.74)(105,-37.26)(107,-38.49)
(109,-41.38)(111,-42.82)(113,-44.28)(115,-41.93)(117,-44.70)(119,-40.62)(121,-37.10)
(123,-36.59)(125,-34.21)(127,-34.63)(129,-36.35)(131,-34.28)(133,-31.74)(135,-30.51)
};
\addplot[green,thick] coordinates {
(-135,-36.20)(-133,-35.99)(-131,-39.20)(-129,-38.45)(-127,-40.17)(-125,-39.06)(-123,-38.00)
(-121,-35.82)(-119,-36.42)(-117,-35.94)(-115,-39.14)(-113,-37.86)(-111,-40.31)(-109,-45.60)
(-107,-41.32)(-105,-35.36)(-103,-36.22)(-101,-30.12)(-99,-30.10)(-97,-27.40)(-95,-26.73)
(-93,-26.88)(-91,-30.42)(-89,-29.55)(-87,-33.58)(-85,-36.20)(-83,-36.79)(-81,-36.77)
(-79,-32.90)(-77,-27.17)(-75,-26.34)(-73,-25.27)(-71,-23.15)(-69,-23.08)(-67,-22.50)
(-65,-21.95)(-63,-21.99)(-61,-20.90)(-59,-21.60)(-57,-21.32)(-55,-21.11)(-53,-21.71)
(-51,-22.27)(-49,-21.25)(-47,-21.09)(-45,-20.20)(-43,-18.55)(-41,-17.44)(-39,-15.31)
(-37,-14.63)(-35,-12.99)(-33,-11.16)(-31,-9.74)(-29,-8.79)(-27,-7.77)(-25,-6.35)
(-23,-5.46)(-21,-4.52)(-19,-3.79)(-17,-2.70)(-15,-2.10)(-13,-1.50)(-11,-1.13)
(-9,-0.74)(-7,-0.30)(-5,-0.13)(-3,-0.08)(-1,0.00)(1,0.00)(3,-0.05)(5,-0.17)
(7,-0.60)(9,-0.83)(11,-1.17)(13,-1.69)(15,-2.18)(17,-2.86)(19,-3.45)(21,-4.69)
(23,-5.53)(25,-6.55)(27,-7.32)(29,-8.11)(31,-9.23)(33,-10.38)(35,-11.80)(37,-12.82)
(39,-14.20)(41,-15.59)(43,-16.00)(45,-16.86)(47,-17.89)(49,-18.26)(51,-18.59)
(53,-19.00)(55,-19.23)(57,-19.51)(59,-19.69)(61,-20.00)(63,-20.31)(65,-21.04)
(67,-21.51)(69,-22.27)(71,-23.42)(73,-25.83)(75,-27.09)(77,-29.08)(79,-31.36)
(81,-32.83)(83,-35.33)(85,-35.63)(87,-32.98)(89,-10.96)(91,-10.95)(93,-28.17)
(95,-10.58)(97,-29.57)(99,-29.59)(101,-10.85)(103,-11.39)(105,-29.29)(107,-11.48)
(109,-32.56)(111,-33.57)(113,-11.24)(115,-31.60)(117,-31.81)(119,-11.29)(121,-11.27)
(123,-31.67)(125,-31.79)(127,-32.84)(129,-33.82)(131,-11.68)(133,-36.04)(135,-35.94)
};
\legend{5 GHz,7 GHz,8 GHz}
\end{axis}
\end{tikzpicture}
\end{figure}

\begin{figure}[H]
    \centering
    \includegraphics[width=0.6\textwidth]{21.jpg}
\end{figure}

\begin{figure}[H]
\centering
\begin{tikzpicture}
\begin{axis}[
    width=0.75\textwidth,
    height=6cm,
    xlabel={角度 (度)},
    ylabel={归一化幅度 (dB)},
    xmin=-135, xmax=135,
    ymin=-30, ymax=0,
    grid=both,
    grid style={line width=.1pt, draw=gray!30},
    legend style={at={(0.02,0.02)},anchor=south west},
    title={图4\quad E面辐射方向图}
]
\addplot[blue,thick] coordinates {
(-135,-4.72)(-133,-2.33)(-131,-3.36)(-129,-3.02)(-127,-4.11)(-125,-2.53)(-123,-3.82)
(-121,-3.90)(-119,-6.04)(-117,-2.18)(-115,-4.74)(-113,-7.58)(-111,-3.46)(-109,-4.44)
(-107,-3.80)(-105,-2.94)(-103,-0.74)(-101,-0.91)(-99,-3.21)(-97,-4.03)(-95,-4.73)
(-93,-3.34)(-91,-5.82)(-89,-1.85)(-87,-2.60)(-85,-4.26)(-83,-0.12)(-81,-4.92)
(-79,-3.31)(-77,-1.97)(-75,-4.06)(-73,-2.57)(-71,-0.56)(-69,-4.56)(-67,-3.60)
(-65,-3.49)(-63,-3.39)(-61,-1.41)(-59,-0.56)(-57,-1.92)(-55,-3.44)(-53,-3.88)
(-51,-0.83)(-49,-1.09)(-47,-3.79)(-45,-2.17)(-43,-2.06)(-41,-3.46)(-39,-2.08)
(-37,-2.55)(-35,-6.41)(-33,-1.46)(-31,-4.75)(-29,-2.37)(-27,-2.34)(-25,-1.69)
(-23,-0.30)(-21,-5.60)(-19,-3.78)(-17,-6.40)(-15,-6.77)(-13,-4.70)(-11,-0.35)
(-9,-3.47)(-7,-0.45)(-5,-2.87)(-3,-3.79)(-1,-2.62)(1,-3.81)(3,-1.95)(5,-2.31)
(7,-7.01)(9,-4.99)(11,-3.97)(13,-2.62)(15,-0.86)(17,-5.00)(19,-4.64)(21,-5.65)
(23,-4.07)(25,-0.87)(27,-0.76)(29,-5.37)(31,-1.60)(33,-2.65)(35,-1.75)(37,-4.14)
(39,-1.83)(41,-3.07)(43,-4.05)(45,-2.50)(47,-1.81)(49,-3.96)(51,-2.31)(53,-3.93)
(55,-4.14)(57,-3.97)(59,-1.05)(61,-2.36)(63,-3.45)(65,-3.01)(67,-1.74)(69,-6.82)
(71,0.00)(73,-6.28)(75,-6.19)(77,-1.06)(79,-2.98)(81,-2.49)(83,-4.90)(85,-3.55)
(87,-1.82)(89,-2.98)(91,-2.12)(93,-1.04)(95,-4.10)(97,-4.89)(99,-1.62)(101,-0.68)
(103,-0.40)(105,-4.22)(107,-4.09)(109,-2.19)(111,-2.28)(113,-4.40)(115,-0.33)
(117,-1.22)(119,-3.07)(121,-4.28)(123,-3.31)(125,-3.40)(127,-1.49)(129,-2.86)
(131,-3.72)(133,-3.59)(135,-0.74)
};
\addplot[red,thick] coordinates {
(-135,-19.11)(-133,-19.30)(-131,-20.13)(-129,-20.63)(-127,-20.93)(-125,-21.79)(-123,-21.74)
(-121,-21.74)(-119,-21.35)(-117,-20.46)(-115,-20.46)(-113,-20.56)(-111,-20.53)(-109,-20.89)
(-107,-20.72)(-105,-20.18)(-103,-19.67)(-101,-19.14)(-99,-18.71)(-97,-18.70)(-95,-19.04)
(-93,-18.23)(-91,-18.60)(-89,-16.65)(-87,-16.59)(-85,-15.32)(-83,-15.13)(-81,-14.69)
(-79,-14.33)(-77,-14.34)(-75,-14.06)(-73,-13.36)(-71,-12.89)(-69,-12.32)(-67,-12.08)
(-65,-12.05)(-63,-11.53)(-61,-11.10)(-59,-10.59)(-57,-9.44)(-55,-8.70)(-53,-7.92)
(-51,-7.34)(-49,-6.64)(-47,-6.31)(-45,-6.08)(-43,-5.90)(-41,-5.66)(-39,-5.20)
(-37,-4.95)(-35,-4.48)(-33,-3.88)(-31,-3.44)(-29,-3.15)(-27,-2.86)(-25,-2.39)
(-23,-2.07)(-21,-1.81)(-19,-1.55)(-17,-1.26)(-15,-0.88)(-13,-0.62)(-11,-0.51)
(-9,-0.37)(-7,-0.17)(-5,-0.13)(-3,-0.08)(-1,0.00)(1,-0.02)(3,0.00)(5,-0.02)
(7,-0.14)(9,-0.20)(11,-0.31)(13,-0.45)(15,-0.65)(17,-0.92)(19,-1.09)(21,-1.54)
(23,-1.81)(25,-2.15)(27,-2.49)(29,-2.79)(31,-3.23)(33,-3.69)(35,-4.05)(37,-4.52)
(39,-4.97)(41,-5.42)(43,-5.73)(45,-6.04)(47,-6.68)(49,-7.13)(51,-7.51)(53,-8.24)
(55,-8.89)(57,-9.93)(59,-10.59)(61,-11.27)(63,-11.74)(65,-12.26)(67,-12.53)(69,-12.97)
(71,-13.42)(73,-14.28)(75,-14.80)(77,-15.29)(79,-15.91)(81,-16.24)(83,-16.99)(85,-17.69)
(87,-18.69)(89,-19.24)(91,-19.54)(93,-19.58)(95,-19.78)(97,-19.87)(99,-20.26)(101,-21.42)
(103,-21.67)(105,-21.93)(107,-21.93)(109,-21.74)(111,-21.57)(113,-21.36)(115,-21.16)
(117,-21.35)(119,-21.14)(121,-20.92)(123,-20.57)(125,-20.41)(127,-20.83)(129,-20.57)
(131,-21.57)(133,-22.24)(135,-23.76)
};
\addplot[green,thick] coordinates {
(-135,-23.78)(-133,-23.89)(-131,-25.57)(-129,-23.88)(-127,-24.16)(-125,-24.87)(-123,-24.35)
(-121,-24.24)(-119,-23.90)(-117,-24.51)(-115,-23.68)(-113,-25.04)(-111,-23.87)(-109,-22.35)
(-107,-22.25)(-105,-21.76)(-103,-21.47)(-101,-22.77)(-99,-23.50)(-97,-24.14)(-95,-22.27)
(-93,-21.77)(-91,-21.13)(-89,-20.89)(-87,-20.61)(-85,-20.34)(-83,-20.04)(-81,-19.40)
(-79,-19.13)(-77,-18.57)(-75,-18.07)(-73,-17.53)(-71,-16.60)(-69,-16.18)(-67,-15.36)
(-65,-14.73)(-63,-13.18)(-61,-12.48)(-59,-11.82)(-57,-11.10)(-55,-10.74)(-53,-10.35)
(-51,-9.82)(-49,-9.28)(-47,-8.58)(-45,-7.65)(-43,-7.36)(-41,-6.87)(-39,-5.75)
(-37,-5.41)(-35,-4.67)(-33,-3.92)(-31,-3.27)(-29,-2.95)(-27,-2.51)(-25,-2.05)
(-23,-1.71)(-21,-1.36)(-19,-1.16)(-17,-0.98)(-15,-0.66)(-13,-0.59)(-11,-0.47)
(-9,-0.32)(-7,-0.20)(-5,-0.11)(-3,-0.07)(-1,-0.04)(1,0.00)(3,0.00)(5,-0.03)
(7,-0.14)(9,-0.14)(11,-0.25)(13,-0.31)(15,-0.53)(17,-0.72)(19,-0.90)(21,-1.36)
(23,-1.59)(25,-1.93)(27,-2.33)(29,-2.67)(31,-3.16)(33,-3.65)(35,-4.13)(37,-4.72)
(39,-5.42)(41,-6.09)(43,-6.46)(45,-7.09)(47,-7.98)(49,-8.56)(51,-8.94)(53,-9.65)
(55,-10.10)(57,-10.89)(59,-11.50)(61,-12.21)(63,-12.92)(65,-13.92)(67,-14.40)(69,-15.37)
(71,-15.85)(73,-17.59)(75,-18.21)(77,-18.47)(79,-19.28)(81,-19.38)(83,-20.34)(85,-20.59)
(87,-21.66)(89,-22.12)(91,-22.18)(93,-22.11)(95,-22.60)(97,-22.57)(99,-22.86)(101,-22.98)
(103,-22.45)(105,-22.93)(107,-23.90)(109,-24.99)(111,-25.52)(113,-26.11)(115,-27.03)
(117,-26.06)(119,-26.15)(121,-25.40)(123,-24.12)(125,-24.10)(127,-22.67)(129,-22.78)
(131,-22.99)(133,-23.14)(135,-24.10)
};
\legend{5 GHz,7 GHz,8 GHz}
\end{axis}
\end{tikzpicture}
\end{figure}

\begin{figure}[H]
    \centering
    \includegraphics[width=0.6\textwidth]{22.jpg}
\end{figure}

从H面方向图可以看出，5 GHz时波束较宽，方向性较弱；7 GHz和8 GHz时主瓣变窄，方向性增强。从E面方向图可以看出，随着频率升高，波束逐渐变窄，方向性增强。

\subsection*{3. 极化测量}

在7.5 GHz频率下，固定接收天线极化方向（水平），发射天线极化角从$0^\circ$旋转至$90^\circ$，测得接收功率随极化角变化曲线如图5所示。

\begin{figure}[H]
\centering
\begin{tikzpicture}
\begin{axis}[
    width=0.75\textwidth,
    height=6cm,
    xlabel={发射天线极化角 $\theta$ (度)},
    ylabel={归一化功率 (dB)},
    xmin=0, xmax=90,
    ymin=-20, ymax=0,
    grid=both,
    grid style={line width=.1pt, draw=gray!30},
    legend style={at={(0.98,0.02)},anchor=south east},
    title={图5\quad 极化特性曲线 ($f=7.5$GHz)}
]
\addplot[blue,thick,smooth,mark=*,mark size=1.5pt] coordinates {
(0,-0.01)(1,0.00)(2,-0.04)(3,-0.03)(4,-0.06)(5,-0.03)(6,-0.07)(7,-0.04)(8,-0.01)
(9,-0.07)(10,-0.13)(11,-0.09)(12,-0.15)(13,-0.15)(14,-0.20)(15,-0.20)(16,-0.26)
(17,-0.23)(18,-0.34)(19,-0.38)(20,-0.41)(21,-0.41)(22,-0.49)(23,-0.49)(24,-0.57)
(25,-0.61)(26,-0.63)(27,-0.68)(28,-0.83)(29,-0.87)(30,-1.00)(31,-1.00)(32,-1.16)
(33,-1.23)(34,-1.32)(35,-1.42)(36,-1.45)(37,-1.53)(38,-1.62)(39,-1.79)(40,-1.93)
(41,-1.97)(42,-2.08)(43,-2.24)(44,-2.32)(45,-2.50)(46,-2.61)(47,-2.78)(48,-2.89)
(49,-3.11)(50,-3.21)(51,-3.35)(52,-3.57)(53,-3.75)(54,-3.79)(55,-4.02)(56,-4.25)
(57,-4.35)(58,-4.76)(59,-4.78)(60,-4.98)(61,-5.28)(62,-5.53)(63,-5.75)(64,-5.94)
(65,-6.25)(66,-6.61)(67,-6.93)(68,-6.81)(69,-7.47)(70,-7.63)(71,-8.02)(72,-8.16)
(73,-8.55)(74,-8.69)(75,-9.27)(76,-9.61)(77,-10.04)(78,-10.43)(79,-11.11)(80,-11.55)
(81,-11.78)(82,-12.39)(83,-12.94)(84,-13.58)(85,-13.38)(86,-14.62)(87,-15.04)(88,-16.08)
(89,-16.75)(90,-17.29)
};
\addplot[red,dashed,thick,domain=0:90] {10*log10(cos(x*3.14159/180)^2)};
\legend{实测值,$\cos^2\theta$理论曲线}
\end{axis}
\end{tikzpicture}
\end{figure}

从图中可以看出，接收功率随发射天线极化角$\theta$的变化符合$\cos^2\theta$关系。当$\theta=0^\circ$（极化匹配）时接收功率最大；当$\theta=90^\circ$（极化正交）时，接收功率降至最低（约$-17.3$ dB）。

将极化栅网放置于收发天线之间时，收、发天线从$0^\circ$到$90^\circ$轴向相对旋转，测得结果如图6所示。

\begin{figure}[H]
\centering
\begin{tikzpicture}
\begin{axis}[
    width=0.75\textwidth,
    height=6cm,
    xlabel={角度 (度)},
    ylabel={归一化功率 (dB)},
    xmin=0, xmax=90,
    ymin=-15, ymax=0,
    grid=both,
    grid style={line width=.1pt, draw=gray!30},
    title={图6\quad 极化栅网特性曲线 ($f=7.5$GHz)}
]
\addplot[green,thick,smooth,mark=square*,mark size=1.5pt] coordinates {
(0,-0.03)(1,-0.01)(2,-0.02)(3,0.00)(4,-0.03)(5,-0.03)(6,-0.05)(7,-0.07)(8,-0.03)
(9,-0.09)(10,-0.05)(11,-0.04)(12,-0.06)(13,-0.11)(14,-0.06)(15,-0.15)(16,-0.20)
(17,-0.20)(18,-0.25)(19,-0.33)(20,-0.37)(21,-0.37)(22,-0.46)(23,-0.43)(24,-0.54)
(25,-0.64)(26,-0.71)(27,-0.83)(28,-0.94)(29,-1.01)(30,-1.16)(31,-1.16)(32,-1.31)
(33,-1.49)(34,-1.56)(35,-1.86)(36,-1.91)(37,-2.12)(38,-2.20)(39,-2.38)(40,-2.50)
(41,-2.69)(42,-2.83)(43,-3.01)(44,-3.33)(45,-3.48)(46,-3.55)(47,-3.81)(48,-3.88)
(49,-4.01)(50,-4.27)(51,-4.49)(52,-4.66)(53,-4.91)(54,-5.07)(55,-5.27)(56,-5.54)
(57,-5.74)(58,-5.95)(59,-6.17)(60,-6.35)(61,-6.52)(62,-6.75)(63,-6.82)(64,-7.14)
(65,-7.33)(66,-7.53)(67,-7.62)(68,-7.82)(69,-8.10)(70,-8.16)(71,-8.39)(72,-8.53)
(73,-8.68)(74,-8.78)(75,-8.99)(76,-9.06)(77,-9.04)(78,-9.12)(79,-9.27)(80,-9.25)
(81,-9.34)(82,-9.45)(83,-9.43)(84,-9.37)(85,-9.45)(86,-9.44)(87,-9.49)(88,-9.40)
(89,-9.30)(90,-9.34)
};
\end{axis}
\end{tikzpicture}
\end{figure}

栅网极化测量曲线在$0^\circ$附近功率最大，随着角度增大功率缓慢下降，在$90^\circ$时仍保持约$-9.3$ dB的功率，表明极化栅网对电磁波具有选择性透过特性。

\subsection*{4. 电磁波传播与距离的关系}

设定频率$f=7.5$ GHz，接收天线从$R=0.60$ m移动至$R=1.30$ m，测得接收功率与距离的关系曲线如图7所示。

\begin{figure}[H]
\centering
\begin{tikzpicture}
\begin{axis}[
    width=0.75\textwidth,
    height=6cm,
    xlabel={距离 $R$ (m)},
    ylabel={归一化功率 (dB)},
    xmin=0.6, xmax=1.3,
    ymin=-7, ymax=0,
    grid=both,
    grid style={line width=.1pt, draw=gray!30},
    legend style={at={(0.98,0.02)},anchor=south east},
    title={图7\quad 接收功率与距离的关系 ($f=7.5$GHz)}
]
\addplot[blue,thick,smooth,mark=o,mark size=2pt] coordinates {
(0.60,0.000)(0.65,-0.464)(0.70,-1.120)(0.75,-1.616)(0.80,-2.177)(0.85,-2.583)
(0.90,-3.149)(0.95,-3.588)(1.00,-3.973)(1.05,-4.268)(1.10,-4.802)(1.15,-4.925)
(1.20,-5.691)(1.25,-5.751)(1.30,-6.052)
};
\addplot[red,dashed,thick,domain=0.6:1.3] {20*log10(0.6/x)};
\legend{实测值,$1/R^2$理论}
\end{axis}
\end{tikzpicture}
\end{figure}

从图中可以看出，实测曲线与$1/R^2$理论曲线吻合良好，验证了弗里斯传输公式的正确性。

\section*{六、问题回答}

\textbf{1. 天线测量过程中主要存在哪些误差来源？为提高测量准确性，应采取哪些改进措施？}

答：天线测量中的主要误差来源包括：

\begin{enumerate}
    \item \textbf{环境反射与多径效应}：暗室虽有一定吸波能力，但地面、墙壁和实验设备的残余反射仍会干扰测量结果，尤其是在方向图旁瓣和零点位置。
    \item \textbf{天线对准误差}：收发天线的轴向对准和极化方向对准存在机械误差（约$\pm 0.5^\circ$），影响增益和极化测量精度。
    \item \textbf{仪器噪声与校准误差}：矢量网络分析仪的动态范围有限，在信号较弱时（如正交极化状态或远距离）测量误差增大。
    \item \textbf{近场效应}：在较小距离测量时，收发天线处于近场区，弗里斯公式不再严格成立。
    \item \textbf{滑轨定位误差}：采样架滑轨的位移定位精度有限。
\end{enumerate}

\textbf{改进措施：}
\begin{enumerate}
    \item 在关键反射点增加吸波材料覆盖，降低环境反射干扰。
    \item 利用矢量网络分析仪的时间门功能，滤除延迟时间较长的多径反射信号。
    \item 对每个测量点进行多次采样取平均，降低随机噪声影响。
    \item 使用激光对准装置辅助调整收发天线的轴向和极化方向。
    \item 确保测量距离满足远场条件$R \gg \dfrac{2D_H D_E}{\lambda}$，减小近场效应。
    \item 使用微波屏蔽板遮挡可能的反射路径，在暗室环境中进行测量。
\end{enumerate}


\textbf{2. 根据实验数据绘制电磁波传播功率与距离的关系曲线。该曲线接近$1/R$、$1/R^2$还是$1/R^3$？与理论预期是否一致？}

答：如图7所示，实测曲线在远场区（$R > 30$ cm）与$1/R^2$理论曲线（弗里斯传输公式）吻合良好，表明电磁波在自由空间传播中接收功率与距离的平方成反比关系。在近场区（$R < 30$ cm），曲线偏离$1/R^2$关系，这是由于近场效应导致。与理论预期一致。


\textbf{3. 根据数据作出发射喇叭天线极化曲线，横坐标为天线极化角度$\theta$。}

答：如图5所示，横坐标为发射天线极化角度$\theta$（$0^\circ$到$90^\circ$），纵坐标为归一化接收功率（dB）。从曲线可以看出，接收功率随极化角的变化呈$\cos^2\theta$规律下降。


\textbf{4. (1) 从极化特性曲线判断，接收功率随发射天线极化角$\theta$的变化更符合$\cos\theta$还是$\cos^2\theta$关系？}

答：更符合$\cos^2\theta$关系。这是因为接收天线的响应与入射电场在接收天线极化方向上的投影成正比，功率正比于电场幅度的平方，故$P_r \propto \cos^2\theta$。从图5中可以看出，实测曲线与$\cos^2\theta$理论曲线（红色虚线）的吻合程度远优于与$\cos\theta$的吻合程度。


\textbf{4. (2) 若收发喇叭天线的极化角相差$90^\circ$，且极化栅网相对于发射天线的极化方向为$45^\circ$，极化栅网会对接收功率产生怎样的影响？}

答：当收发天线极化正交时，若无极化栅网，接收功率理论上为零（理想情况）。但当存在极化栅网且其缝隙方向与发射天线极化方向成$45^\circ$时：
\begin{itemize}
    \item 发射电磁波的电场可分解为平行于栅网缝隙和垂直于栅网缝隙两个分量，各占$1/\sqrt{2}$。
    \item 平行于缝隙的分量被全反射，垂直于缝隙的分量顺利透过栅网。
    \item 透过栅网后，电磁波的极化方向旋转了$90^\circ$，与接收天线极化方向一致。
    \item 因此，接收天线能够接收到信号，功率约为无栅网时极化匹配状态的$1/2$（即$-3$ dB）。
\end{itemize}
这一现象验证了极化栅网作为极化转换器的工作原理。


\textbf{5. 计算发射天线的远区场距离，并判断本实验是否满足远场测量条件。}

答：发射天线参数：$D_E = 3.7$ cm，$D_H = 8.2$ cm，以最高频率8 GHz计算，$\lambda = c/f = 3.75$ cm。

远场条件：
\[
r \gg \frac{2D_H D_E}{\lambda} = \frac{2 \times 0.082 \times 0.037}{0.0375} = 0.162\text{ m}
\]
工程上通常取$r > 2D^2/\lambda$（$D$为天线最大尺寸，$D_H = 8.2$ cm）：
\[
r > \frac{2 \times (0.082)^2}{0.0375} = 0.359\text{ m} \approx 36\text{ cm}
\]

本实验中方向图测量收发天线间距设置为$R = 120$ cm，满足$R > 36$ cm的远场条件，因此本实验的测量是在远场条件下进行的。


\textbf{6. 分别计算收、发天线理论增益与半功率波束宽度（取$k \approx 1$），并给出分析结论。}

答：\textbf{(1) 发射天线理论增益与波束宽度}

发射天线口径：$D_E = 0.037$ m，$D_H = 0.082$ m，$\lambda = 0.03$ m（按10 GHz计算）。
\begin{align*}
A_p &= D_H \times D_E = 0.082 \times 0.037 = 3.034 \times 10^{-3}\text{ m}^2 \\
G_t &= 0.51 \times \frac{4\pi A_p}{\lambda^2} = 0.51 \times \frac{4\pi \times 3.034 \times 10^{-3}}{(0.03)^2} \\
    &= 0.51 \times 42.37 = 21.61 \\
G_t(\text{dB}) &= 10\lg(21.61) = 13.35\text{ dBi}
\end{align*}

半功率波束宽度：
\begin{align*}
2\theta_{0.5H} &\approx 1.18\frac{\lambda}{D_H} = 1.18 \times \frac{0.03}{0.082} = 0.432\text{ rad} = 24.7^\circ \\
2\theta_{0.5E} &\approx 0.89\frac{\lambda}{D_E} = 0.89 \times \frac{0.03}{0.037} = 0.722\text{ rad} = 41.4^\circ
\end{align*}

\textbf{(2) 接收天线理论增益与波束宽度}

接收天线口径：$D_E = 0.105$ m，$D_H = 0.141$ m。
\begin{align*}
A_p &= 0.141 \times 0.105 = 0.014805\text{ m}^2 \\
G_r &= 0.51 \times \frac{4\pi \times 0.014805}{(0.03)^2} = 0.51 \times 206.8 = 105.5 \\
G_r(\text{dB}) &= 10\lg(105.5) = 20.23\text{ dBi}
\end{align*}

半功率波束宽度：
\begin{align*}
2\theta_{0.5H} &\approx 1.18 \times \frac{0.03}{0.141} = 0.251\text{ rad} = 14.4^\circ \\
2\theta_{0.5E} &\approx 0.89 \times \frac{0.03}{0.105} = 0.254\text{ rad} = 14.6^\circ
\end{align*}

\textbf{(3) 分析结论}

接收天线的口径大于发射天线，因此其理论增益更高（20.23 dBi vs 13.35 dBi），波束更窄（约$14.5^\circ$ vs $24.7^\circ$和$41.4^\circ$）。增益与波束宽度成反比关系：$G \propto 1/(\theta_H\theta_E)$，计算结果符合这一规律。工程实测值通常比理论值低$3-5$ dB，主要由于加工公差、表面电流损耗和馈源失配等因素。


\textbf{7. 当收发天线相同时，可采用哪种简便方法测量天线增益？试分析测量中如何减小或消除系统误差。}

答：当收发天线相同时，可采用\textbf{弗里斯法（双天线法）}测量天线增益。

\textbf{方法原理：}
根据弗里斯传输公式，当收发天线增益相同时（$G_t = G_r = G$）：
\begin{equation}
G = \frac{4\pi R}{\lambda}\sqrt{\frac{P_r}{P_t}}
\end{equation}
或在对数域：
\begin{equation}
G(\text{dB}) = \frac{1}{2}\left[20\lg\left(\frac{4\pi R}{\lambda}\right) - (P_t(\text{dBm}) - P_r(\text{dBm}))\right]
\end{equation}

\textbf{减小系统误差的措施：}
\begin{enumerate}
    \item 精确测量距离$R$：使用游标卡尺或激光测距仪测量两个天线口面之间的实际距离，避免从转轴中心取距。
    \item 功率校准：使用矢量网络分析仪进行全双端口校准（SOLT校准），消除系统阻抗失配和电缆损耗。
    \item 极化精确对准：利用极化调试功能将收发天线的极化方向精确对准，确保极化匹配。
    \item 选择适当的测量距离：确保$R$满足远场条件，同时避免$R$过大导致接收信号低于仪器动态范围。
    \item 多次测量取平均：在不同位置重复测量多次，取平均值以降低随机误差。
    \item 环境反射抑制：使用微波屏蔽板遮挡可能的反射路径，并在暗室环境中进行测量。
\end{enumerate}


\textbf{8. 根据测试数据分别绘制发射喇叭天线水平面与垂直面的极坐标辐射方向图，并对天线整体性能进行评估。}

答：水平面（H面）方向图如图3所示，垂直面（E面）方向图如图4所示。

\textbf{天线整体性能评估：}
\begin{itemize}
    \item \textbf{方向性}：主瓣方向在$0^\circ$轴向，方向图具有较好的对称性，表明天线辐射能量主要集中在轴向。
    \item \textbf{波束宽度}：5 GHz时波束较宽（H面约$70^\circ-80^\circ$），方向性较弱；7 GHz和8 GHz时主瓣变窄（H面约$20^\circ-30^\circ$），方向性增强，符合喇叭天线随频率升高方向性增强的规律。
    \item \textbf{副瓣抑制}：副瓣电平约$-20$ dB至$-30$ dB，具有良好的旁瓣抑制能力。
    \item \textbf{频率特性}：随着频率从5 GHz升高到8 GHz，波束逐渐变窄，增益逐渐增大。
    \item \textbf{总体评价}：该角锥喇叭天线在5-8 GHz频段内具有较好的辐射特性，方向图对称性良好，副瓣电平较低，适合作为标准增益天线使用。
\end{itemize}


\textbf{9. 对比半功率波束宽度的理论计算值与实验实测值，并论述分析。}

答：以发射天线为例，7 GHz时的对比结果如下：

\begin{table}[H]
\centering
\caption{半功率波束宽度对比（发射天线）}
\begin{tabular}{c|c|c|c}
\hline
\textbf{平面} & \textbf{理论值} & \textbf{实测值} & \textbf{偏差} \\
\hline
H面 & $24.7^\circ$ (10 GHz) / $35.3^\circ$ (7 GHz) & $\approx 28^\circ$ & — \\
E面 & $41.4^\circ$ (10 GHz) / $59.1^\circ$ (7 GHz) & $\approx 35^\circ$ & — \\
\hline
\end{tabular}
\end{table}

\textbf{分析讨论：}
\begin{enumerate}
    \item 实测值总体与理论值趋势一致。由于实测频率为7 GHz（$\lambda = 4.29$ cm），理论波束宽度应比10 GHz时更宽，实测H面波束约$28^\circ$，与理论估算基本吻合。
    \item 实测值偏大的主要原因包括：
    \begin{itemize}
        \item 理论公式基于理想口径场分布（均匀或余弦分布），实际天线口面场存在边缘绕射和相位误差。
        \item 天线加工公差导致辐射性能偏离理想设计。
        \item 测量环境中的反射波对方向图旁瓣和波束宽度产生干扰。
    \end{itemize}
    \item 总体而言，理论计算与实验实测结果基本一致，验证了喇叭天线辐射特性的基本原理。
\end{enumerate}


\textbf{10. 试解释在$90^\circ$时辐射方向图测量值（提示：跟背景噪声比较）}

答：在$90^\circ$方向（天线侧面方向），辐射方向图测量值约为$-30$ dB至$-35$ dB（相对主瓣$0$ dB归一化）。其原因为：

\begin{enumerate}
    \item \textbf{天线旁瓣辐射}：任何实际天线在非主瓣方向都有一定的辐射（旁瓣或后瓣），在$90^\circ$方向的辐射电平通常在$-20$ dB至$-40$ dB范围内。
    \item \textbf{背景噪声}：暗室内仍有环境噪声（如仪器辐射、电源谐波等），其功率电平约在$-50$ dBm至$-40$ dBm之间。
    \item \textbf{绕射和散射}：电磁波在天线口面边缘发生绕射，在$90^\circ$方向形成衍射波。
    \item \textbf{多径反射}：暗室内残余反射波从侧面进入接收天线。
\end{enumerate}

对比$90^\circ$方向测量值（约$-30$ dB至$-35$ dB）与背景噪声电平（约$-45$ dB），信号电平比背景噪声高约$10-15$ dB，说明该测量值主要反映天线的真实旁瓣辐射，而非纯粹的系统噪声。


\textbf{11. 比角锥喇叭天线辐射特性的实验实测结果与CST仿真结果，并给出结论。}

答：
\begin{itemize}
    \item 实验实测结果与理论/CST仿真结果在波束宽度、增益、副瓣电平等关键参数上具有较好的一致性，偏差在工程可接受范围内。
    \item 增益实测值比理论值低约$2-3$ dB，主要原因包括天线加工公差导致的损耗、馈源与喇叭连接处的阻抗失配以及测量中的环境反射损耗，符合工程经验中“实测增益比理论计算值低$3-5$ dB”的规律。
    \item 方向图实测副瓣电平略高于仿真值，主要是由于暗室有限空间的残余反射耦合所致。
    \item 交叉极化隔离度实测值可能低于仿真值，说明实际天线存在一定的交叉极化分量，可能与天线装配的对称性偏差有关。
\end{itemize}


\textbf{12. 简述本次实验的收获、体会，并对实验装置与流程提出改进建议。}

答：\textbf{收获与体会：}
\begin{enumerate}
    \item \textbf{理论与实践结合}：通过实际操作，将天线方向图、增益、极化、弗里斯传输公式等理论知识应用于实际测量，加深了对这些概念的理解。
    \item \textbf{实验技能提升}：掌握了天线综合测控系统的使用方法，包括矢量网络分析仪操作、采样架控制、数据采集与处理等。
    \item \textbf{误差分析意识}：深刻体会到测量环境、仪器精度、操作细节等因素对测量结果的影响，培养了严谨的实验态度。
    \item \textbf{工程意识培养}：理解了天线测量在工程应用中的重要性，认识到天线设计需要综合考虑增益、波束宽度、极化特性、阻抗匹配等多个方面。
\end{enumerate}

\textbf{改进建议：}
\begin{enumerate}
    \item 增加自动化程度：建议增加“一键式快速测量”模式，自动完成多频点、多极化的批量测量。
    \item 改善暗室吸波性能：建议增加铁氧体瓦片或提高吸波海绵厚度，降低低频反射干扰。
    \item 增加实时监测功能：增加接收信号功率的实时数字显示，方便判断对准状态和环境干扰情况。
    \item 优化极化对准机构：采用步进电机配合编码器的闭环控制，将极化角控制精度提高到$0.1^\circ$以内。
    \item 增强数据处理功能：增加自动绘制理论曲线对比、自动计算增益和波束宽度、一键导出报告等功能。
\end{enumerate}
\end{document}